% ============================================================
% section6_angular.tex - Раздел 6: Теоретическая модель: программная реализация
% ============================================================

\section{Теоретическая модель: программная реализация}

В предыдущем разделе мы разобрали теоретические основы работы проволочных поляризаторов и матричный формализм для системы из двух решеток. Теперь пришло время перевести эту электродинамическую модель на язык Python. Мы напишем модуль, реализующий расчеты по статьям Blanco et al. (1986) и Kaltenecker et al. (2019).

\subsection{Реализация формул Бланко: от теории к коду}

Для расчёта теоретического пропускания мы создадим модуль \texttt{theoretical.py}. В нём нам потребуется вычислить функции $A_1$, $A_2$, $f_a$, $f_b$, $f_c$ и $f_d$, которые, в свою очередь, зависят от суммы бесконечных рядов. На практике достаточно ограничиться первыми $N=15$ членами ряда, чтобы получить очень высокую точность.

Давай сопоставим ключевые формулы из статьи Blanco с их реализацией на Python:

\textbf{1. Функция $C_m$} (Формула 12 из Blanco):
\begin{equation*}
C_m = \left[ m^2 - \left(\frac{p}{\lambda}\right)^2 \right]^{1/2}
\end{equation*}
Поскольку для высоких частот (когда шаг решетки $p$ превышает длину волны $\lambda$) подкоренное выражение может стать отрицательным, $C_m$ становится комплексным числом. Поэтому в коде мы явно добавляем мнимую единицу \texttt{+ 0j}:
\begin{minted}[linenos=false]{python}
def compute_C(m, p_over_lambda):
    val = m**2 - p_over_lambda**2
    return np.sqrt(val + 0j)
\end{minted}

\textbf{2. Параметр $A_1$} (Формула 10 из Blanco):
В статье формула выглядит так:
\begin{equation*}
A_1 = 1 + \frac{1}{2}\left(\frac{\pi d}{\lambda}\right)^2 \left( \ln\frac{p}{\pi d} + \frac{3}{4} \right) + \frac{1}{2}\left(\frac{\pi d}{\lambda}\right)^2 \sum_{m=1}^\infty \left( \frac{1}{C_m} - \frac{1}{m} \right)
\end{equation*}
В коде мы разбиваем её на три слагаемых (\texttt{term1}, \texttt{term2}, \texttt{term3}) и используем защищенный логарифм \texttt{safe\_log} для предотвращения математических ошибок (например, деления на ноль или логарифма от отрицательного числа), которые могут возникнуть из-за ошибок округления:
\begin{minted}[linenos=false]{python}
def safe_log(x):
    return np.log(max(x, EPS))
\end{minted}

\textbf{3. Параметр $A_2$} (Формула 11 из Blanco):
По аналогии вычисляется и второй параметр:
\begin{equation*}
A_2 = 1 + \frac{1}{2}\left(\frac{\pi d}{\lambda}\right)^2 \left( \frac{11}{4} - \ln\frac{p}{\pi d} \right) + \frac{1}{24}\left(\frac{\pi d}{\lambda}\right)^2 - \left(\frac{\pi d}{\lambda}\right)^2 \sum_{m=1}^\infty \left( m - \frac{1}{2m}\left(\frac{p}{\lambda}\right)^2 - C_m \right)
\end{equation*}
Как и в предыдущем пункте, программная реализация разбивает это длинное выражение на отдельные слагаемые.

\textbf{4. Эквивалентные проводимости $f_a$ и $f_b$} (Формулы 8 и 9 из Blanco):
Эти функции определяют мнимые части импедансов:
\begin{align*}
f_a(p, d, \lambda) &= \frac{p}{2\lambda}\left(\frac{\pi d}{p}\right)^2 \frac{1}{A_2} \\
f_b(p, d, \lambda) &= \frac{2\lambda}{p}\left(\frac{p}{\pi d}\right)^2 A_1 - \frac{p}{4\lambda}\left(\frac{\pi d}{p}\right)^2 \frac{1}{A_2}
\end{align*}
\note{Обрати внимание на опечатку в оригинальной статье Blanco (1986): там формула для $f_b$ ошибочно пронумерована как (19), хотя должна быть (9). В более свежей работе Kaltenecker (2019) эта формула приведена корректно в системе уравнений (5).}

\subsection{Вычисление амплитудных коэффициентов пропускания}

Следующим шагом является расчет комплексных амплитудных коэффициентов $T_\perp$ и $T_\parallel$ по формулам (3) и (13) из Blanco. 

\begin{minted}{python}
def compute_t_perp(p_over_lambda, d_over_p, N=15, tol=1e-8):
    # ... (вычисление fa, fb) ...
    Za = -1j / fa
    Zb = -1j / fb
    Z1 = (Za**2 + Za*Zb + Za**2*Zb) / (2*Za + Za**2 + Zb + Za*Zb + EPS)
    Z2 = Za / (1.0 + Za)
    denom = (1.0 + Z1) * (Zb + Z2)
    return 2.0 * Z1 * Z2 / denom if np.abs(denom) > EPS else 0.0
\end{minted}
Здесь мы вводим малую константу \texttt{EPS = 1e-12}, чтобы избежать деления на ноль. Это критически важно, так как в некоторых особых точках (например, при целочисленных значениях $p/\lambda$, которые соответствуют резонансам Вуда) импедансы могут обращаться в бесконечность или ноль. Для таких особых точек в функции предусмотрены специальные ветвления (возврат значения с допуском \texttt{tol}), реализующие формулы (20) и (21) из статьи Blanco.

\subsection{Матрица Джонса для двух поляризаторов}

Система из двух поляризаторов, где первый повернут на угол $\theta$, описывается уравнением (1) из статьи Kaltenecker. В Python матричные операции элегантно реализуются через оператор \texttt{@} (матричное умножение) из библиотеки NumPy:

\begin{minted}{python}
def transmission_two_polarizers(theta, p_over_lambda, d_over_p, N=15):
    t_perp = compute_t_perp(p_over_lambda, d_over_p, N)
    t_par = compute_t_par(p_over_lambda, d_over_p, N)

    c = np.cos(theta)
    s = np.sin(theta)

    # Матрицы вращения (R и R_inv) и матрицы поляризатора (P)
    R = np.array([[c, -s], [s, c]])
    R_inv = np.array([[c, s], [-s, c]])
    P = np.array([[t_perp, 0], [0, t_par]])

    # Полная матрица Джонса: M = P * R(theta) * P * R(-theta)
    M = P @ R @ P @ R_inv

    # Входное поле поляризовано вдоль оси X
    E_in = np.array([1.0, 0.0])
    E_out = M @ E_in

    # Энергетический коэффициент пропускания (мощность)
    T = np.abs(E_out[0])**2 + np.abs(E_out[1])**2
    return np.clip(T, 0.0, 1.0) # Защита от выхода за пределы [0, 1]
\end{minted}

\subsection{Модифицированная модель для реального эксперимента}

Идеальная теория Бланко предполагает абсолютно ровные, идеально проводящие бесконечные проволоки. На практике решетка имеет дефекты: неровности намотки, конечное сопротивление металла, рассеяние на шероховатостях.

Чтобы приблизить теоретическую кривую к экспериментальным данным, мы вводим \textbf{модифицированную модель} в функции \texttt{transmission\_db\_modified}. Она содержит два подгоночных параметра:
\begin{itemize}
    \item \texttt{period\_scale} — масштабный множитель для периода. Из-за нерегулярности намотки эффективный шаг решетки $p_{eff}$ может отличаться от паспортного значения.
    \item \texttt{loss\_factor} — коэффициент частотно-зависимых потерь. Он моделирует экспоненциальное затухание, растущее с частотой, обусловленное омическими потерями и рассеянием.
\end{itemize}

\begin{minted}{python}
def transmission_db_modified(theta_deg, p, d, freq_THz, period_scale, loss_factor, N=15):
    # ...
    p_eff = p * period_scale
    T_linear = transmission_two_polarizers(theta_rad, p_eff/lambda_m, d/p_eff, N)
    
    # Добавляем эмпирические потери
    T_linear_mod = T_linear * np.exp(-loss_factor * freq_THz)
    return 10 * np.log10(np.maximum(T_linear_mod, 1e-12))
\end{minted}

\subsection{Полный листинг модуля}

Для продолжения работы полный код модуля \texttt{theoretical.py} вынесен в \textbf{\textit{Приложение~\ref{app:theoretical}}}. Этот модуль станет нашим вычислительным «ядром» для расчёта теоретических кривых. Скопируй его в файл \texttt{theoretical.py} в корне твоего проекта.

\subsection{Итоги раздела}
Мы успешно перевели электродинамическую модель Бланко и матричный формализм Кальтенекера в надежный Python-код. Особое внимание было уделено защите от математических сингулярностей (деление на ноль, отрицательные значения под логарифмом) и добавлению эмпирических параметров (\texttt{period\_scale} и \texttt{loss\_factor}) для приближения теории к реальности. Наше расчетное ядро \texttt{theoretical.py} готово к работе! В следующем разделе мы объединим эти расчеты с экспериментальными данными и построим наглядную угловую зависимость.
